% Generated by scripts/render_publication_curvature.py from saved JSON; do not edit.
The probe evaluates the frozen legacy curvature surrogate $J_B$ at the saved IBM01 endpoint of the first curvature stage, in float32 and without any optimization or legalization. At that point $J_B=0.86905$ and $\lVert\nabla J_B\rVert_2=8.09\times10^{-2}\,\mu\mathrm{m}^{-1}$. A Lanczos run capped at 14 Hessian-vector products (HVPs) of the automatic-differentiation (AD) operator returned Rayleigh quotient $-0.442458$ with residual norm $3.49\times10^{-7}$. The historical cascade recorded the smallest eigenvalue $-0.442458$ at the same endpoint; the two values differ by $1.49\times10^{-7}$.

Let $x$ be the saved coordinate vector, $g_{\mathrm{AD}}(x)$ its AD gradient, and $H_{\mathrm{AD}}v$ the AD HVP for a unit vector $v$ over the movable coordinates. The central differences are
\[d_{\mathrm{FD}}(\epsilon)=\frac{J_B(x+\epsilon v)-J_B(x-\epsilon v)}{2\epsilon},\qquad k_{\mathrm{FD}}(\epsilon)=v^\top\frac{g_{\mathrm{AD}}(x+\epsilon v)-g_{\mathrm{AD}}(x-\epsilon v)}{2\epsilon}.\]
Positions and $\epsilon$ are in micrometers, so slopes are in $\mu\mathrm{m}^{-1}$ and curvatures in $\mu\mathrm{m}^{-2}$; the surrogate itself is dimensionless.

\begin{table}[htbp]
\centering\footnotesize
\caption{Float32 central differences along two seeded random unit directions at the saved IBM01 endpoint, all movable coordinates included. The AD columns are the slope and curvature of the branch active at $x$; the FD columns sample both sides of nearby branch boundaries. Slopes are in $\mu\mathrm{m}^{-1}$ and curvatures in $\mu\mathrm{m}^{-2}$. The directions are independent of the Lanczos vector.}
\label{tab:ibm01-curvature-diagnostic}
\begin{tabular}{@{}rrrrrr@{}}
\toprule
Dir. & $\epsilon$ ($\mu$m) & $g_{\mathrm{AD}}^\top v$ & $d_{\mathrm{FD}}$ & $v^\top H_{\mathrm{AD}}v$ & $k_{\mathrm{FD}}$ \\
\midrule
1 & $1.00\times10^{-3}$ & $-1.33\times10^{-4}$ & $-5.36\times10^{-4}$ & $3.44\times10^{-3}$ & $-2.88\times10^{-1}$ \\
1 & $3.00\times10^{-3}$ & $-1.33\times10^{-4}$ & $-5.66\times10^{-4}$ & $3.44\times10^{-3}$ & $-1.01\times10^{-1}$ \\
1 & $1.00\times10^{-2}$ & $-1.33\times10^{-4}$ & $-5.33\times10^{-4}$ & $3.44\times10^{-3}$ & $-1.28\times10^{-2}$ \\
1 & $3.00\times10^{-2}$ & $-1.33\times10^{-4}$ & $-5.00\times10^{-4}$ & $3.44\times10^{-3}$ & $9.50\times10^{-3}$ \\
1 & $1.00\times10^{-1}$ & $-1.33\times10^{-4}$ & $-5.35\times10^{-4}$ & $3.44\times10^{-3}$ & $1.82\times10^{-2}$ \\
2 & $1.00\times10^{-3}$ & $-7.68\times10^{-4}$ & $-3.28\times10^{-4}$ & $4.93\times10^{-4}$ & $-3.12\times10^{-2}$ \\
2 & $3.00\times10^{-3}$ & $-7.68\times10^{-4}$ & $-2.09\times10^{-4}$ & $4.93\times10^{-4}$ & $2.07\times10^{-3}$ \\
2 & $1.00\times10^{-2}$ & $-7.68\times10^{-4}$ & $-1.88\times10^{-4}$ & $4.93\times10^{-4}$ & $1.49\times10^{-2}$ \\
2 & $3.00\times10^{-2}$ & $-7.68\times10^{-4}$ & $-2.17\times10^{-4}$ & $4.93\times10^{-4}$ & $1.71\times10^{-2}$ \\
2 & $1.00\times10^{-1}$ & $-7.68\times10^{-4}$ & $-3.15\times10^{-4}$ & $4.93\times10^{-4}$ & $2.09\times10^{-2}$ \\
\bottomrule
\end{tabular}
\end{table}

For the full vectors, with $h_{\mathrm{FD}}=[g_{\mathrm{AD}}(x+\epsilon v)-g_{\mathrm{AD}}(x-\epsilon v)]/(2\epsilon)$, the relative difference $\lVert h_{\mathrm{FD}}-H_{\mathrm{AD}}v\rVert_2/\lVert H_{\mathrm{AD}}v\rVert_2$ ranges from 7.96 to 776 over the ten rows, and the cosine between the two vectors from $-1.36\times10^{-1}$ to $6.74\times10^{-2}$. At the endpoint, 21 coordinates sit at or beyond a canvas clamp. For direction~1 at $\epsilon=10^{-3}\,\mu$m the forward and backward cost slopes are $-8.34\times10^{-4}$ and $-2.38\times10^{-4}\,\mu\mathrm{m}^{-1}$, so the two sides already lie on different branches.

% Follow-up source: writeup/data/publication/curvature_followup_ibm01.json.
\paragraph{Precision and branch follow-up.} A separate saved probe restricts both directions to the 2{,}233 coordinates whose canvas clearance exceeds $0.1\,\mu$m in both precisions, so no perturbation can cross a canvas clamp. All forward floating operations run in float64; stored float32 constants are promoted without regeneration. With $P$ selecting these coordinates, the comparison uses $PH_{\mathrm{AD}}Pv$ and differences of $Pg_{\mathrm{AD}}$. The frozen composite is also split into its three weighted primitives, which reproduce the composite value, slope and curvature to $10^{-14}$. Table~\ref{tab:ibm01-curvature-followup} reports the gradient jump $\lVert P[g(x+\epsilon v)-g(x-\epsilon v)]\rVert_2$ in $\mu\mathrm{m}^{-1}$.

\begin{table}[htbp]
\centering\footnotesize
\caption{Guarded float64 follow-up at the saved IBM01 endpoint: gradient jump across $\pm\epsilon$ (in $\mu\mathrm{m}^{-1}$) and the relative difference between the finite-difference and AD Hessian-vector products at $\epsilon=10^{-5}\,\mu$m. A smooth term has a jump proportional to $\epsilon$; a jump that does not shrink with $\epsilon$ marks a branch boundary.}
\label{tab:ibm01-curvature-followup}
\begin{tabular}{@{}llrrrr@{}}
\toprule
Dir. & Term & $\epsilon=10^{-7}$ & $\epsilon=10^{-5}$ & $\epsilon=10^{-3}$ & Rel.\ HVP diff. \\
\midrule
1 & Wirelength & $1.51\times10^{-10}$ & $1.51\times10^{-8}$ & $1.51\times10^{-6}$ & $3.91\times10^{-9}$ \\
1 & $\tfrac12\times$density & $6.01\times10^{-3}$ & $2.20\times10^{-2}$ & $3.85\times10^{-2}$ & $4.52\times10^{5}$ \\
1 & $\tfrac12\times$RUDY & $2.50\times10^{-3}$ & $2.61\times10^{-3}$ & $2.61\times10^{-3}$ & $5.17\times10^{3}$ \\
1 & Composite $J_B$ & $6.51\times10^{-3}$ & $2.22\times10^{-2}$ & $3.86\times10^{-2}$ & $4.33\times10^{4}$ \\
\addlinespace
2 & Wirelength & $1.54\times10^{-10}$ & $1.54\times10^{-8}$ & $1.54\times10^{-6}$ & $3.45\times10^{-9}$ \\
2 & $\tfrac12\times$density & $6.01\times10^{-3}$ & $2.15\times10^{-2}$ & $3.84\times10^{-2}$ & $4.11\times10^{5}$ \\
2 & $\tfrac12\times$RUDY & $2.50\times10^{-3}$ & $2.61\times10^{-3}$ & $2.17\times10^{-2}$ & $3.98\times10^{3}$ \\
2 & Composite $J_B$ & $6.51\times10^{-3}$ & $2.16\times10^{-2}$ & $4.42\times10^{-2}$ & $3.28\times10^{4}$ \\
\bottomrule
\end{tabular}
\end{table}

Inside the guarded region the wirelength gradient changes in proportion to $\epsilon$, and its finite-difference HVP agrees with the AD product to $3.91\times10^{-9}$ and $3.45\times10^{-9}$ at $\epsilon=10^{-5}\,\mu$m. The density and RUDY gradients jump by amounts that persist down to $\epsilon=10^{-7}\,\mu$m, where the $\pm\epsilon$ samples already differ in 311 recorded min/max/clamp comparisons of the composite in each direction. At $\epsilon=10^{-3}\,\mu$m the composite jump norms are $3.86\times10^{-2}$ and $4.42\times10^{-2}$, and the relative HVP differences are 752.48 and 670.50.
